\part{Introducción al Cálculo diferencial}

\null\vfill
\begin{Huge}\begin{center}
IV. Cálculo Diferencial	
\end{center}\end{Huge}

\vspace{1.5cm}

\begin{figure}[H]
	\centering
	\includegraphics[width=.9\textwidth]{imagenes/part4.png}	
\end{figure}
\par
\vfill


\chapter{Introducción al Cálculo diferencial}
\begin{tikzpicture}
	\fill [left color=red!50, right color=teal!50] (0,0) rectangle (6.5,.2);
	\fill [left color=teal!50, right color=blue!50] (6.5,0) rectangle (11.5,.2);
	\end{tikzpicture}

\vspace{10mm}


Los temas de introducción al cálculo diferencial: ``Funciones'', ``Continuidad y Límites'', ``La derivada''; pueden verse en 

\begin{Large}
\hspace{2cm} \textcolor{NavyBlue}{https://github.com/igvaori/calculo-infinitesimal}
\end{Large}

\vspace{5mm} Si se desea incluir en el temario de `Precálculo' algo de probabilidad y estadística, se puede acudir a:

\begin{Large}
\hspace{2cm} \textcolor{NavyBlue}{https://github.com/igvaori/Estadistica}
\end{Large}

\vspace{5mm} Para otros contenidos, incluso de Física, se puede echar un vistazo general a:

\begin{Large}
\hspace{2cm}\textcolor{NavyBlue}{https://github.com/igvaori}	
\end{Large}


















\begin{comment}

%%%%%%%%%%%%%%%%%%%%%%%%%%%%%%%%%%%. SECCIONES
\chapter{texto}
\begin{tikzpicture}
	\fill [left color=red!50, right color=teal!50] (0,0) rectangle (6.5,.2);
	\fill [left color=teal!50, right color=blue!50] (6.5,0) rectangle (11.5,.2);
	\end{tikzpicture}

\vspace{1cm}
\section{texto}
\begin{tikzpicture}
	\fill [left color=red!50, right color=teal!50] (0,0) rectangle (3.5,.1);
	\fill [left color=teal!50, right color=blue!50] (3.5,0) rectangle (7.5,.1);
	\end{tikzpicture}
\vspace{0.5cm}

\subsection{texto}
\begin{tikzpicture}
	\fill [left color=red!50, right color=teal!50] (0,0) rectangle (3.5,.01);
	\fill [left color=teal!50, right color=blue!50] (3.5,0) rectangle (7.5,.01);
	\end{tikzpicture}
\vspace{0.5cm}


%%%%%%%%%%%%%%%%%%%%%%%%%%%%%%%%%%%. \begin{ ------>. 
detsacado;  cuadro-naranja;  cuadro-gris;  miejercicio (solución extensa);  mipropuesto (solución corta y fuera del cuadro)

%%%%%%%%%%%%%%%%%%%%%%%%%%%%%%%%%%%. CURIOSIDAD
\vspace{1cm}
\color{ForestGreen!80}
\rule{250pt}{0.2pt}
Texto
\vspace{-8mm}
\begin{flushright}
\rule{250pt}{0.2pt}		
\end{flushright}	
\color{black}
\end{comment}